\section{PDDL planning}
% TODO: fill after PDDL experiments. (PDDL use and integration = 20%.)
% Content plan (the exam asks to identify WHERE and WHEN PDDL helps):
%  - Where: PDDL as a plan-library member serving the go_to intention;
%    the domain (move actions on a directed tile graph) and how arrow
%    tiles are encoded by omitting edges.
%  - When: deterministic, fully-believed subproblems; why the outer loop
%    is a poor fit (per-frame world changes, online solver latency).
%  - Problem generation from beliefs (reachable region, size bound).
%  - Measured comparison: BFS plan vs PDDL plan (latency, plan length,
%    failure rate) from plannerCalls/plannerFailures logs.
%  - Discussion of problem complexity, as required.
%  - Future work / design alternative (DECISION, deliberate scope): a richer
%    PDDL model could plan whole collect-and-deliver sequences (pickup,
%    putdown, capacity, mission constraints). We deliberately keep PDDL at
%    the go_to level because the environment is highly dynamic (long plans
%    go stale within seconds) and the online solver adds latency; task
%    selection stays with the BDI strategy, while PDDL provides a meaningful
%    but bounded planning component. Crate-pushing is likewise out of scope
%    (the agent assumes a static known map + partial sensing).
